\pagenumbering{arabic}
\section{INTRODUCCIÓN}
%La introducción es una guía para el lector, donde se presenta lo que encontrará al leer el documento.
%De esta sección (incluida la introducción en el conteo total de páginas) comienza el conteo de páginas para un máximo de 60 páginas en el documento.
El pronóstico de demanda condiciona directamente las políticas de reposición de inventario: cuánto pedir, cuándo hacerlo y cuánto mantener como stock de seguridad. Cuando los tiempos de aprovisionamiento son elevados o variables, subestimar la demanda provoca quiebres de stock, mientras que sobrestimarla inmoviliza capital, adicional otros elementos como el efecto látigo (\textit{bullwhip effect}) amplifica estos errores a lo largo de la cadena de suministro.

La literatura ofrece varias familias de modelos para este problema: estadísticos clásicos, de machine learning, de aprendizaje profundo y fundacionales de series de tiempo que operan por inferencia \textit{zero-shot}. Sin embargo, su desempeño depende de la forma de la serie y del horizonte de pronóstico, y sin una evaluación controlada sobre un mismo conjunto de datos la elección termina dependiendo más de la familiaridad tecnológica que de la evidencia experimental.

Este trabajo aborda la siguiente pregunta de investigación: \textit{¿Qué familia de modelos de pronóstico ofrece el mejor desempeño para la demanda de inventarios de una empresa a partir de sus datos históricos, y en qué perfiles de demanda esa superioridad puede verificarse empíricamente mediante backtesting retrospectivo sobre una ventana temporal de validación real?}

Para responderla se desarrolló PRED (\textit{...}), un framework experimental que evalua los modelos para cada SKU buscando el modelo que mejor rendimiento brinde, esto aplicado sobre el caso experimental de el inventario de una clinica.

El documento se organiza así: la sección II presenta la oportunidad, el problema y los objetivos; la sección III el marco teórico y el estado del arte; la sección IV el análisis del problema y los requerimientos; la sección V el diseño de la solución; la sección VI su desarrollo; la sección VII los resultados; y la sección VIII las conclusiones y el trabajo futuro.